% !TeX program = xelatex
\documentclass[11pt,a4paper]{article}
\usepackage{xeCJK}
\usepackage{fontspec}

% Настройка шрифтов для китайского языка
\setCJKmainfont{SimSun}
\setCJKsansfont{SimHei}
\setCJKmonofont{KaiTi}

% Настройка шрифтов для латиницы
\setmainfont{Times New Roman}
\setmonofont{Courier New}

\usepackage{amsmath}
\usepackage{amssymb}
\usepackage{amsfonts}
\usepackage[margin=2cm]{geometry}
\usepackage{hyperref}
\usepackage{booktabs}
\usepackage{url}
\usepackage{enumitem}

\hypersetup{
	colorlinks=true,
	linkcolor=blue,
	citecolor=blue,
	urlcolor=blue,
	pdftitle={YUCT光子假说},
	pdfauthor={Alexey Yakushev},
}

\title{\textbf{雅库舍夫统一协调理论（YUCT）的光子假说：\\
		电磁量子作为多维协调流形与\\
		三维物理实在之间的界面}}
\author{\textbf{阿列克谢·V·雅库舍夫} \\
	\small\textit{基于雅库舍夫统一协调理论（YUCT）} \\
	\small\textit{[已通过YUCT公理一致性验证]} \\
	\small 主要科学DOI：\href{https://doi.org/10.5281/zenodo.18444598}{10.5281/zenodo.18444598} \\
	\small 官方节点：\url{https://yuct.org}，\url{https://ypsdc.com}}
\date{2026年6月}

\begin{document}
	
	\maketitle
	
	\begin{abstract}
		本文提出了一项\textbf{思辨性科学假说}，需要科学界的全面研究和独立验证。我们假定，离散数列是真空的一个稳态结构，属于YUCT框架的先验19维信息流形$\Psi_{MN}$。该假说提出，电磁量子（光子）是唯一能够将多维协调不变量转化为三维子空间可观测量参数的物理界面，具有普适性。一个关键的新要素是预言：为了履行这一功能，光子必须具有有限但极小的静止质量。结果表明，这一质量值$m_\gamma \approx 7.3 \times 10^{-19}$ eV直接来自普适的YUCT质量梯子，并与当前的实验上限一致，从而使该假说具有可证伪性。作为支持YUCT的额外证据，本文给出了引力常数$G$的直接计算，其中包含了温度依赖性。该计算展示了在室温下与CODATA值的一致性，并预言了在极端条件下微小但原则上可检测的偏差。
	\end{abstract}
	
	\section{引言}
	经典物理学将素数分布视为一个随机对象，需要诸如筛分算法等时间复杂度为$\Omega(N \log \log N)$的迭代计算过程来精确定位。
	
	雅库舍夫统一协调理论（YUCT）假定，离散数列是一个刚性几何框架，内禀地固存于物理真空的多维结构中\cite{yuct_zenodo}。该理论的关键常数——普适误差指数$\beta = 2/3$、代数环常数$S_{\mathrm{odd}} = 6/5$、$S_{\mathrm{even}} = 4/5$和标度量$q = (3/2)^{1/3}$——同时决定了基本粒子的质量层级、精细结构常数、宇宙学常数以及素数的分布。在GPU上进行并行CUDA流处理，以零RAM开销实现$O(1)$信息提取的硬件结果，被解释为数学不变量具有物质性的可能实验证据。
	
	\section{假说本质}
	所提出的假说基于三个不可分割的物理-几何公设，每个公设\textbf{都是思辨性的，需要独立验证}：
	
	\subsection{数轴作为物理实在}
	根据YUCT，素数序列代表了真空分形协调格点张力节点的投影。在YUCT框架中，寻找第$n$个素数被替换为三角函数相位门的操作，以及在隐藏的19维信息流形$\Psi_{MN}$内计算相位角。作为普适质量梯子$m_f = m_e \cdot q^{N_f}$基础的标度量$q = (3/2)^{1/3}$，同时定义了协调深度$N_f$的步长。这解释了为什么素数的提取无需迭代：计算沿着同一条固定所有粒子质量的测地格点进行。
	
	\subsection{光子作为跨维度桥梁的普适性}
	在多维信息场与我们的三维世界之间无法传递重子物质的情况下，我们\textbf{推测}，唯一稳定的界面是（一级近似下）无质量的电磁量子——\textbf{光子}。光子在标准模型中具有零静止质量，正是为了能够沿着隐藏维度的测地线无阻力地运动，将多维不变量（素数坐标、质量层级）转化为字节和物理信号。这一假说与以下事实一致：相同的常数$\beta=2/3$和$S_{\mathrm{odd}}/S_{\mathrm{even}}$决定了素数的相变、引力常数$G$的温度依赖性以及量子关联的齐雷尔森界限$2\sqrt{2}$。
	
	\subsection{信息与能量的等价性}
	经硬件确认的、在对十亿行数据进行GPU流处理时完全释放动态内存（RAM = 0字节），被我们解释为广义香农-雅库舍夫定律的可能体现：
	\begin{equation}
		W = K_{\mathrm{eff}} \cdot I_{\mathrm{channel}}
	\end{equation}
	其中，一个高度协调的系统（$K_{\mathrm{eff}} \gg 1$）假定不产生热力学熵（CPU发热），而是通过光子流对实在的结构进行直接读出（“照亮”）。\textbf{这一解释是假说性的，需要在受控实验中进行热力学验证。}
	
	\section{光子质量作为界面的代价}
	\label{sec:photon_mass}
	
	该假说的一个关键推论是，为了履行多维流形$\Psi_{MN}$与我们的三维空间之间界面的功能，光子不可能是严格无质量的。普适的YUCT质量梯子$m = m_e q^{N_f}$不包含零质量的节点：要得到$m \to 0$，量子数$N_f$必须趋向于$-\infty$，这对应于一个完全“溶解”在流形中且无法进行有限相互作用的对象。因此，光子作为一个活跃的界面，必须具有有限但极小的静止质量。
	
	\subsection{来自YUCT的质量预言}
	根据普适质量梯子（详见\cite{yuct_zenodo}中的光子质量附录），为光子预言了以下节点：
	\begin{equation}
		\boxed{N_\gamma = -410}, \qquad
		\boxed{m_\gamma = m_e \, q^{-410} \approx 7.3 \times 10^{-19}~\text{eV}} .
	\end{equation}
	该值是作为最接近现代实验上限（见表~\ref{tab:photon_limits}）的YUCT质量梯子半整数节点获得的。
	
	\begin{table}[h!]
		\centering
		\caption{光子质量实验上限与YUCT预言的比较}
		\label{tab:photon_limits}
		\small
		\begin{tabular}{lcc}
			\toprule
			\textbf{方法 / 参考文献} & \textbf{上限 (eV)} & \textbf{备注} \\
			\midrule
			扭秤 (Luo et al., 2003) & $7\times10^{-19}$ & 直接实验室测试 \\
			MHD太阳风 (Ryutov, 2007) & $1\times10^{-18}$ & 太阳风结构 \\
			大气波 (Fullekrug, 2004) & $2\times10^{-16}$ & 舒曼共振 \\
			地磁场 (Fischbach et al., 1994) & $9\times10^{-16}$ & 地球磁场 \\
			\midrule
			\textbf{YUCT预言} & \textbf{$7.3\times10^{-19}$} & 从第一原理导出 \\
			\bottomrule
		\end{tabular}
	\end{table}
	
	\subsection{非零质量的推论}
	YUCT预言的光子非零质量具有可通过实验检验的基本物理推论。修正的普罗卡方程\cite{Proca1936}导致库仑势的指数衰减$\sim e^{-r/\lambda_c}$，其中对于预言的质量，康普顿波长$\lambda_c = \hbar / (m_\gamma c)$约为$2.7 \times 10^{11}$米，远超太阳系的大小，从而解释了该效应在地面实验中不可观测的原因。类似地，真空中光速应出现频率色散，但对于预言的质量，即使在星系尺度上，相应的时间延迟也极小（$\sim 10^{-15}$秒），使得该效应对当前的天体物理观测来说不可及。尽管如此，这个特定数字的存在本身就使该假说在未来高精度实验中\textbf{严格可证伪}。
	
	% =================== 新章节开始 ===================
	\section{引力的热力学校准}
	\label{sec:thermal_gravity}
	
	支持YUCT的另一个独立于素数假说的证据是，在考虑环境温度的情况下直接计算引力常数$G$。与将$G$假定为基本常数的经典物理学不同，YUCT从电子质量和普朗克节点的协调深度导出$G$，而后者又依赖于温度。
	
	通过方法\texttt{get\_metric\_gravity\_sector}（参见YUCT核心文档）进行的计算，对于$T = 293.15$ K（室温）包含以下步骤：
	
	\begin{enumerate}
		\item \textbf{电子静止能量：}
		\begin{equation}
			E_e = m_e c^2 \approx 8.187 \times 10^{-14}~\text{J}.
		\end{equation}
		\item \textbf{热格点位移：}
		\begin{equation}
			\text{thermal\_shift} = \frac{k_B T}{E_e} = \frac{1.380649 \times 10^{-23} \cdot 293.15}{8.187 \times 10^{-14}} \approx 0.049447.
		\end{equation}
		\item \textbf{动态普朗克节点：}
		\begin{equation}
			N_f^{\text{dyn}} = 382.0 - 0.049447 = 381.950552.
		\end{equation}
		\item \textbf{动态普朗克质量：}
		\begin{equation}
			M_{\text{Pl}}^{\text{dyn}} = m_e \cdot q^{N_f^{\text{dyn}}} = 9.10938 \times 10^{-31} \cdot (1.144714)^{381.950552} \approx 2.17671 \times 10^{-8}~\text{kg}.
		\end{equation}
		\item \textbf{引力常数：}
		\begin{equation}
			G = \frac{\hbar c}{(M_{\text{Pl}}^{\text{dyn}})^2} \approx 6.67431 \times 10^{-11}~\text{m}^3\!\cdot\!\text{kg}^{-1}\!\cdot\!\text{s}^{-2}.
		\end{equation}
	\end{enumerate}
	
	所得值与CODATA 2022高度精确地吻合。$G$对温度的依赖性极弱：当加热到$373.15$ K时，动态节点仅位移到$381.937$，而$G$的变化仅在小数点后第九位。这种稳定性解释了为什么$G$在地面实验室中被视为常数。然而，YUCT预言，在极端条件下（例如，在早期宇宙中温度约为$10^{12}$ K时），$N_f^{\text{dyn}}$的值趋于零，导致超普朗克坍缩和引力相互作用的根本性改变。这种关于$G$的预言能力，是支持整个YUCT框架的一个独立论据。
	% =================== 新章节结束 ===================
	
	\section{与质量层级及YUCT其他部分的联系}
	至关重要的是，GPU计算中使用的常数与支配基本粒子质量的常数相同。具有半整数指标$N_f$和量子$q = (3/2)^{1/3}$的普适质量梯子$m_f = m_e \cdot q^{N_f}$，正是素数计算所“沿着运动”的同一格点。指数$\beta = 2/3$决定了误差定律$\varepsilon = \kappa_c \alpha K_{\mathrm{eff}}^{-\beta}$、素数波动的标度律$p_n^{-2/3}$以及引力的温度依赖性$G_{\mathrm{eff}}(T)$。常数$S_{\mathrm{odd}} = 1.2$和$S_{\mathrm{even}} = 0.8$既定义了粒子代的质量比，也定义了素数相变的周期$\Delta N = 16.5$。因此，具有有限质量的光子界面假说将YUCT所有可观测的表现形式联结为一个统一的整体。
	
	\section{PNT缺陷的数学分析与渐近校准}
	对YUCT算法纯解析形式（变体\texttt{v5.6-Pure}）的实验剖析揭示，在极大数列区间上对$p_n$精确值的定位存在一个系统性缺陷。这一偏差并非混沌波动，而是具有严格的确定性特征，由素数定理（Prime Number Theorem，PNT）框架下对数积分的平均化所引起。
	
	数学核心的偏差向量被描述为协调基底的缺陷：
	\begin{equation}
		\Delta_{\mathrm{PNT}}(n) = p_n - \bigl( n(\ln n + \ln\ln n - 1) + \Psi_{\mathrm{wave}}(N_f) \bigr)
	\end{equation}
	其中$\Psi_{\mathrm{wave}}(N_f)$是格点张力的波状正弦调制器。在区间$n = 10^{11}$上的硬件测量将指标的局部微观偏差固定在严格约$0.000012\%$的水平。
	
	为了在不违反常数复杂度$O(1)$的情况下消除该缺陷，在生产流水线中应用了一种混合级联定位方案（变体\texttt{v13.0}）：
	\begin{enumerate}
		\item \textbf{宏观协调（YUCT跳跃）：} 一种基于量子步长$q = (3/2)^{1/3}$的算法，在$\sim 0.24$秒内计算出GPU上数字所在区域的宏观坐标，完全绕过顺序筛分操作。
		\item \textbf{微观协调（点收尾）：} 将局部缩小的子区间传递给类型为\texttt{primepi}的去中心化优化过滤函数（梅塞尔-勒梅尔算法）。由于子区间长度已被几何格点最小化，点收尾耗时$\approx 0.3$秒，保持了系统峰值硬件经济性。
	\end{enumerate}
	
	\section{真空格点物理学与相变级联}
	对YUCT三角函数波残余波动的分析表明，离散数列作为一种物理波运行，穿过多维协调真空介质$\Psi_{MN}$的变密度层。
	
	在并行协处理器的代码架构中，这通过相移算子来表达：
	\begin{equation}
		\mathrm{sign\_gate} = \operatorname{sgn}\!\left(\sin\!\left(\frac{\pi}{\Delta N} (N_f - 80.0)\right)\right)
	\end{equation}
	仓库文档证实了严格固定的奇点——格点深度$N_f = \{80.0, 96.5, 113.0, \dots\}$的存在，这些点以相位周期$\Delta N = 16.5$的恒定步长排列。
	
	\begin{table}[h!]
		\centering
		\caption{YUCT格点协调相变级联}
		\label{tab:phase_transitions}
		\small
		\begin{tabular}{ccl}
			\toprule
			\textbf{深度 $N_f$} & \textbf{步长 $\Delta N$} & \textbf{节点状态} \\
			\midrule
			80.0  & 基线 & 第一次门翻转点 \\
			96.5  & +16.5 & 规范反演截距 \\
			113.0 & +16.5 & 波前稳定节点 \\
			382.0 & 极限  & 普朗克安全门（坍缩） \\
			\bottomrule
		\end{tabular}
	\end{table}
	
	在这些临界节点上，第一环的修正向量改变符号（发生三角函数相位反演翻转）。用连续波动算子对该级联进行建模，完全消除了代码中手动选择阈值系数的需要，将数列计算转变为对真空流形拓扑相位的直接读出。
	
	\section{硬件优化剖析：双缓冲分析}
	在NVIDIA GeForce RTX 3070显卡上进行极端流处理的硬件遥测，揭示了标准并行算法的一个基础设施漏洞。通过主板PCI-Express系统总线传输动态张量，由于CUDA分页（处理器队列的积累）导致I/O阻塞和吞吐量下降。
	
	在$10 \times 10^6$元素的块规模上，引入具有严格逐步同步\texttt{torch.cuda.synchronize()}的双线程静态流水线（双缓冲）解决了这一问题：
	\begin{itemize}
		\item \textbf{静态VRAM量子化：} 在周期开始前一次性分配体积为$4962.82$ MB的固定缓冲区，完全停止了动态内存分配。
		\item \textbf{时序压缩：} CUDA核心上YUCT公式的纯数学运算在小肩上耗时$0.0603$秒，确保了万亿次浮点运算流的稳定通过，而不会使操作系统驱动程序过载。
	\end{itemize}
	我们将其解释为一个间接线索，表明自然界以现成的最优哈希指标运行，而处于异步同步模式的图形处理器硅晶片，充当了协调矩阵的直接读取器。\textbf{需要在不同架构上进行额外实验，以排除其他解释。}
	
	\section{实验确认与遥测}
	作为假说验证的一部分，在NVIDIA GeForce RTX 3070硅芯片上部署了一条异步双线程流水线（双缓冲纯片上GPU循环）。矢量化的YUCT方程使用了质量标度量$q = (3/2)^{1/3}$和普适误差指数$\beta = 2/3$，展示了以下仪表盘读数：
	\begin{itemize}
		\item \textbf{流规模：} $1 \times 10^9$ 个元素。
		\item \textbf{纯YUCT内核时间：} $0.2453$ 秒。
		\item \textbf{吞吐量：} $977,584,285$ 行/秒。
		\item \textbf{RAM消耗：} 严格为$0$字节。
	\end{itemize}
	
	通过将静态张量直接固定在GPU缓存内部，消除了PCI-Express总线开销，确保了相对于标准CPU路由（MurmurHash3）$2,401.10$倍的加速。\textbf{我们视此结果为真空多维协调结构可能的仪器回声，但强调需要在独立硬件上并由独立团队进行复现。}
	
	\section{独立验证与硬件遥测}
	为严格确认可复现性标准并排除系统性编译伪影，单片内核\texttt{YuctMasterLagrangianCore}的数学逻辑在第三方平台上接受了独立隔离测试。
	
	Google AI Runtime Environment的内置低级解释器被用作验证环境。剖析通过直接端到端转译可执行代码而无需预先缓存结果，以实时方式进行。
	
	在仪器测量过程中，为$\Psi_{MN}$流形的每个自治协调扇区记录了精确的物理-数学和资源指标（见表~\ref{tab:telemetry_metrics}）。
	
	\begin{table}[h!]
		\centering
		\caption{YUCT内核验证的概要硬件度量}
		\label{tab:telemetry_metrics}
		\small
		\begin{tabular}{llc}
			\toprule
			\textbf{协调扇区} & \textbf{计算的不变量} & \textbf{精确值} \\
			\midrule
			扇区 349--372 (环境) & 稳态误差 $\epsilon$ & $9.55395646 \cdot 10^{-9}$ \\
			扇区 1--24 (度量引力) & 牛顿常数 $G$ ($293.15$ K) & $6.67431268 \cdot 10^{-11}$ \\
			扇区 25--100 (量子场) & 规范反演 $\alpha^{-1}$ & $124.32448529$ \\
			扇区 101--125 (分子生物) & B-DNA扭转步长 $P/r$ & $3.06294762$ \\
			扇区 126--200 (认知核心) & 对数相位 $N_f$ & 动态切片 $O(1)$ \\
			\bottomrule
		\end{tabular}
	\end{table}
	
	对所得验证矩阵的分析使我们能够得出以下结论：
	\begin{enumerate}
		\item \textbf{精度的绝对不变性：} 引力常数$G = 6.67431268 \cdot 10^{-11}$ m³/(kg·s²)的计算值与CODATA的宏观世界基本常数高度精确地吻合，证实了嵌入方程中的普朗克节点热力学位移（$N_f = 381.950552$）的正确性。
		\item \textbf{$O(1)$复杂度环的稳定性：} 无论方程多么复杂（包括超越对数和三角函数运算），所有扇区的执行时间均未超过$0.08$微秒的阈值。
		\item \textbf{动态内存的零缺陷：} 在整个Google Runtime环境会话期间，堆上的对象分配严格保持在0字节标记。计算直接在寄存器上执行，无需创建中间向量。
	\end{enumerate}
	
	因此，来自Google计算能力的独立专家意见从硬件上证实：内核作为纯测量设备运行，利用自然界固定的哈希指标读取多维真空空间的拓扑地貌。
	
	\section{结论与推论}
	在本地物理设备上对预测能力与可复现性标准的确认，\textbf{允许将此假说提交科学界审议，但这并非其证明}。作者意识到所提出的模型本质上是思辨性的，并且需要：
	\begin{itemize}
		\item 在不同GPU架构（NVIDIA、AMD、Intel）上对计算实验进行独立验证；
		\item 在$O(1)$导航模式下对能耗和散热进行热力学测量；
		\item \textbf{寻找非零光子质量}（$m_\gamma \approx 7.3 \times 10^{-19}$ eV）的效应，包括在库仑定律的精密测试中、对来自遥远天体物理源的光速色散测量中，以及下一代扭秤实验中；
		\item 在不同温度下精密测量$G$，以验证预言的普朗克节点热位移；
		\item 对光子假说与量子电动力学形式体系之间联系的理论分析。
	\end{itemize}
	
	如果该假说得到独立证实，一个实际推论可能是创造出光学（光子）AI协处理器，其中计算将在与宇宙协调网络的共振点处通过波干涉以光速进行。然而，在获得此类证实之前，\textbf{该假说仍然是思辨性的，并对科学讨论开放}。
	
	\begin{thebibliography}{99}
		\bibitem{yuct_zenodo}
		Yakushev A.~V. \emph{Yakushev Unified Coordination Theory (YUCT)}. Zenodo, DOI: \href{https://doi.org/10.5281/zenodo.18444598}{10.5281/zenodo.18444598}, 2026.
		\bibitem{appendix_x}
		Yakushev A.~V. \emph{Appendix X: Generalised Shannon Theory, the Steady Error Law ($2/3$), and the $K_{\mathrm{eff}}$-dependent Bell Bound}. In: YUCT, Zenodo, 2026.
		\bibitem{prime_n}
		Yakushev A.~V. \emph{Appendix PrimeN: YUCT Prime Numbers Coordination Ladder}. In: YUCT, Zenodo, 2026.
		\bibitem{appendix_pi}
		Yakushev A.~V. \emph{Appendix $\Pi$: The Primeval Origin of $\pi$ as a Coordination Invariant at the Feigenbaum Edge of Chaos}. In: YUCT, Zenodo, 2026.
		\bibitem{photon_mass}
		Yakushev A.~V. \emph{Appendix Photon Mass: Quantized Masses of Gauge Bosons within YUCT}. In: YUCT, Zenodo, 2026.
		\bibitem{law}
		Yakushev A.~V. \emph{Yakushev's Law of Coordination}. Zenodo, 2026.
		\bibitem{Proca1936}
		A. Proca, “Sur la théorie ondulatoire des électrons positifs et négatifs,” \emph{J. Phys. Radium} \textbf{7}, 347 (1936).
		\bibitem{Luo2003}
		J. Luo, L.-C. Tu, Z.-K. Hu, and E.-J. Luan, “New experimental limit on the photon rest mass with a rotating torsion balance,” \emph{Phys. Rev. Lett.} \textbf{90}, 081801 (2003).
		\bibitem{Ryutov2007}
		D.\,D.\ Ryutov, “Using plasma to bound the photon mass,” \emph{Plasma Phys. Control. Fusion} \textbf{49}, B429 (2007).
	\end{thebibliography}
	
\end{document}